\documentclass[12pt]{article}
\usepackage[T1]{fontenc}
\usepackage[utf8]{inputenc}
\usepackage{lmodern}
\usepackage{textcomp}
\usepackage{enumitem}
\usepackage{geometry}
\geometry{margin=1in}
\begin{document}
\begin{center}
Answer Key - Mathematics - Undergraduate Level Assessment 5 \
\small (Answers are ordered exactly as in the question paper.)
\end{center}

\section*{Multiple Choice}
\begin{enumerate}[label=\arabic*.]
\item \textbf{Answer:} B
\\ \textit{Options:} A) commutative semi group; B) abelian group; C) non-abelian group; D) None of these
\\ \textit{Note:} The condition (ab)^-1 = a^-1 b^-1 implies that the operation in the group is commutative, meaning a*b = b*a for all a, b in G, which defines an abelian group.
\item \textbf{Answer:} B
\\ \textit{Options:} A) not anti-symmetric; B) transitive; C) reflexive; D) symmetric
\\ \textit{Note:} The relation \textasciitilde{} is transitive because for every pair (a, b) and (b, c) in the relation, there exists a pair (a, c), specifically, since (1, 2) and (1, 3) do not violate transitivity and the other pairs do not provide a counterexample.
\item \textbf{Answer:} D
\\ \textit{Options:} A) 0; B) 1; C) 0,1; D) 0,4
\\ \textit{Note:} The correct answer is 0 and 4 because these are the values in the finite field Z\_5 that satisfy the polynomial equation $x^5$ + 3 $x^3$ + $x^2$ + 2 x = 0 when evaluated modulo 5.
\item \textbf{Answer:} D
\\ \textit{Options:} A) finite; B) cyclic; C) of order two; D) abelian
\\ \textit{Note:} The condition ($ab)^2$ = $a^2$ $b^2$ for all a, b in G implies that ab * ab = aa * bb, which can be rearranged to show that ab = ba, demonstrating that G is abelian since the operation is commutative for any elements a and b in the group.
\item \textbf{Answer:} B
\\ \textit{Options:} A) 6; B) 12; C) 30; D) 105
\\ \textit{Note:} The maximum possible order of an element in the symmetric group S\_n is determined by the least common multiple of the lengths of the disjoint cycles in its cycle decomposition, and for n=7, the cycle type (3,2,2) produces an order of 12, which is the highest achievable.
\end{enumerate}

\section*{True / False}
\begin{enumerate}[label=\arabic*.]
\setcounter{enumi}{5}
\item \textbf{Answer:} True
\\ \textit{Options:} A) True; B) False
\\ \textit{Note:} The functional equation f(1 + x) = f(x) implies that f is periodic with a period of 1, allowing us to determine that f(15/2) = f(5) = 11.
\item \textbf{Answer:} False
\\ \textit{Options:} A) True; B) False
\\ \textit{Note:} The statement is false because while every permutation can be expressed as a product of disjoint cycles, not every permutation can be expressed as a single disjoint cycle, particularly when considering permutations with multiple independent cycles.
\item \textbf{Answer:} False
\\ \textit{Options:} A) True; B) False
\\ \textit{Note:} The statement is false because the circle intersects the parabola at a maximum of two points, while other curves can intersect the parabola at more than two points, such as lines or other parabolas.
\item \textbf{Answer:} False
\\ \textit{Options:} A) True; B) False
\\ \textit{Note:} The product of (2,3) and (3,5) in the ring Z\_5 x Z\_9 is calculated as (2*3 mod 5, 3*5 mod 9), which results in (1,6), not (3,1), making the statement false.
\item \textbf{Answer:} False
\\ \textit{Options:} A) True; B) False
\\ \textit{Note:} The statement is false because a function is a permutation if and only if it is both one-to-one and onto, and this definition holds true; thus, it is incorrect to claim it is "not always the case."
\end{enumerate}

\section*{Long Form Response}
\begin{enumerate}[label=\arabic*.]
\setcounter{enumi}{10}
\item \textbf{Answer:} Cardinality in set theory refers to the size of a set, specifically the number of elements it contains. When considering the set of all functions from the real numbers R to the set \{0, 1\}, we can observe that each function can be seen as a way to assign either 0 or 1 to every real number. This results in a vast number of possible functions, specifically 2^|R|, where |R| represents the cardinality of the real numbers. Since the cardinality of the continuum, |R|, is uncountably infinite, this leads to the conclusion that the set of all such functions has a cardinality of 2^ℵ₁, which is greater than the cardinality of R itself. Therefore, the set of all functions from R to \{0, 1\} not only demonstrates the concept of cardinality but also highlights why it has the greatest cardinality compared to other sets.
\\ \textit{Note:} correct because it highlights that the set of all functions from R to \{0, 1\} has a cardinality of 2^|R|, which is greater than the cardinality of any other set, due to the uncountably infinite size of R.
\item \textbf{Answer:} The trace of an n \ensuremath{\times} n matrix, denoted as tr(A), is defined as the sum of its diagonal elements. A key property of the trace is that it is linear, meaning tr(A + B) = tr(A) + tr(B) and tr(cA) = c tr(A) for any scalar c. Additionally, the trace of the product of two matrices satisfies the cyclic property: tr(AB) = tr(BA). However, when examining the statement regarding tr($A^2$), it is important to note that while tr($A^2$) is always non- negative (as it is the sum of squares of the eigenvalues of A), the specific assertions about its properties can depend on the context or conditions given. Therefore, the validity of statements regarding tr($A^2$) holds true under certain conditions, particularly when A is symmetric or positive definite, which justifies the assertion that only statement II is valid.
\\ \textit{Note:} correct because it accurately describes the properties of the trace, including linearity and the cyclic property, while also addressing the specific case of tr($A^2$) being non-negative due to its relationship with the eigenvalues of the matrix A.
\end{enumerate}

\vfill
\noindent\textit{End of Answer Key}
\end{document}